% Methodology

\label{sec:methodology}

This section defines the formal reinforcement learning setting and the experimental methodology used in the midway baseline study. The purpose is not to present the final GRPO-control algorithm, which remains part of the final project stage, but to fix the notation, baseline protocol, and implementation setup on which the final comparison will build. In particular, the methodology connects the project motivation in Section~\ref{sec:introduction} with the completed PPO, SAC, and TD3 baseline results reported in Section~\ref{sec:experiments}.

\subsection{Problem Formulation}
\label{sec:methodology-mdp}

The control tasks are modelled as Markov decision processes (MDPs). An MDP is denoted by
\[
\mathcal{M} = (\mathcal{S}, \mathcal{A}, P, r, \gamma, \rho_0),
\]
where $\mathcal{S}$ is the state or observation space, $\mathcal{A}$ is the action space, $P(s' \mid s,a)$ is the transition kernel, $r(s,a)$ is the reward function, $\gamma \in [0,1)$ is the discount factor, and $\rho_0$ is the initial-state distribution. At each time step $t$, the agent observes $s_t$, selects an action $a_t$, receives a reward $r_{t+1}$, and transitions to a new state or observation $s_{t+1}$.

This notation is used for all three project environments. The environments \texttt{cartpole\_swingup} and \texttt{acrobot\_swingup} provide vector-based observations, while \texttt{car\_racing\_continuous} provides image observations. Strictly speaking, the observed input need not always equal the full hidden environment state. However, for the purposes of this midway report, the observed input is treated as the state representation used by the learning algorithm. This is sufficient for defining the baseline training and evaluation protocol, while a more detailed theoretical treatment can be postponed to the final GRPO analysis.

\subsection{Policies, Trajectories, and Returns}
\label{sec:methodology-policies-returns}

A stochastic policy is denoted by $\pi_\theta(a \mid s)$, where $\theta$ are the policy parameters. This notation is used for PPO and SAC. For deterministic policies, as used in TD3, the policy is written as $\mu_\theta(s)$. A policy induces trajectories of the form
\[
\tau = (s_0,a_0,r_1,s_1,a_1,r_2,\ldots).
\]
The discounted return from time $t$ is
\[
G_t^\gamma
=
\sum_{k=0}^{\infty} \gamma^k R_{t+k+1}.
\]
The corresponding state-value and action-value functions are
\[
V^\pi(s)
=
\mathbb{E}_{\pi}
\left[
G_t^\gamma \mid S_t=s
\right],
\]
and
\[
Q^\pi(s,a)
=
\mathbb{E}_{\pi}
\left[
G_t^\gamma \mid S_t=s, A_t=a
\right].
\]
The advantage function is defined as
\[
A^\pi(s,a)
=
Q^\pi(s,a)-V^\pi(s).
\]
This notation is required for describing policy-gradient and actor-critic methods. It is also relevant for the final GRPO-control extension, since GRPO-style methods replace or modify the way the policy update signal is estimated. At the midway stage, however, this report only fixes the notation and baseline protocol; the final GRPO-control variant is not yet implemented.

The reported evaluation metric is distinct from the discounted training return. Following the project specification, learning curves are reported using undiscounted evaluation episode return,
\[
G^{\mathrm{eval}}(\tau)
=
\sum_{t=0}^{H-1} r_{t+1},
\]
where $H$ is the episode horizon. Thus, the figures in this report show observed evaluation return, not training loss or a discounted value estimate.

\subsection{Baseline Algorithms}
\label{sec:methodology-baselines}

The midway study uses PPO, SAC, and TD3 as baseline algorithms. PPO is an on-policy policy-gradient method based on clipped policy updates, SAC is an off-policy maximum-entropy actor-critic method, and TD3 is an off-policy deterministic actor-critic method designed to reduce overestimation errors in continuous control \citep{schulman2017proximal,haarnoja2018soft,fujimoto2018addressing}. In this report, these algorithms are not the proposed methodological contribution. Instead, they define the comparison surface required for the later GRPO-control evaluation.

The baseline choice is useful because the three algorithms represent different design choices in deep reinforcement learning: on-policy versus off-policy learning, stochastic versus deterministic policies, and policy-gradient versus actor-critic optimization. At the midway stage, the purpose is not to establish general algorithmic superiority. The training budgets are limited, and no hyperparameter optimization was performed. The purpose is instead to establish a complete and reproducible baseline matrix against which the final GRPO-control method can later be compared.

\subsection{Environment and Implementation Setup}
\label{sec:methodology-implementation}

The project uses two implementation paths because the environments have different observation types. The vector-control environments \texttt{cartpole\_swingup} and \texttt{acrobot\_swingup} are executed through ObjectRL \citep{baykal2025objectrl}. A project-side environment bridge exposes the selected environments through a Gymnasium-style interface without modifying the external ObjectRL source code. The observations are flattened into one-dimensional \texttt{float32} vectors so that they match the vector-observation workflow expected by the ObjectRL baselines.

The CarRacing environment requires a separate implementation path because it uses image observations. The project therefore implements CNN-based PPO, SAC, and TD3 agents for \texttt{car\_racing\_continuous}. The image observations are converted from height--width--channel format to channel--height--width format and scaled to floating point values,
\[
(96,96,3) \rightarrow (3,96,96).
\]
This makes the observations compatible with convolutional policy and critic networks. The CarRacing runs were executed on Google Colab with CUDA and copied back into the local repository, after which they were processed together with the vector-control runs.

This split is a practical implementation decision. CarRacing is not excluded from the midway study; it is included in the same algorithm--environment--seed matrix as the two vector-control tasks. The important methodological point is that all results are exported into a common report-facing result format, enabling a unified validation and analysis pipeline.

\subsection{Training and Evaluation Protocol}
\label{sec:methodology-protocol}

The midway baseline matrix evaluates three algorithms, three environments, and five random seeds:
\[
3 \times 3 \times 5 = 45
\]
algorithm--environment--seed combinations. The seeds are $0,1,2,3,4$. The final report-facing notebook validates that all 45 combinations are present and that the result set contains 900 evaluation rows.

The vector-control environments are evaluated over a midway budget of 20,000 training steps, while CarRacing is evaluated over 10,000 training steps. Policies are tested at regular evaluation intervals, and the reported metric is undiscounted evaluation episode return. For each algorithm $\alpha$, environment $e$, seed $z$, training step $n$, and evaluation episode $m$, let
$G^{\mathrm{eval}}_{\alpha,e,z,n,m}$ denote the observed evaluation return. The learning curve can then be understood as the empirical mean
\[
\bar{J}_{\alpha,e}(n)
=
\frac{1}{|\mathcal{Z}|M}
\sum_{z \in \mathcal{Z}}
\sum_{m=1}^{M}
G^{\mathrm{eval}}_{\alpha,e,z,n,m},
\]
where $\mathcal{Z}=\{0,1,2,3,4\}$ is the seed set and $M$ is the number of evaluation episodes per evaluation point.

The result files store the algorithm, environment, seed, training step, evaluation episode, evaluation return, completion status, and implementation metadata. This enables the notebook to validate file coverage, detect missing combinations, aggregate final returns, rank algorithms within each environment, and display the learning curves from a single consistent result structure.

\subsection{Preparation for the Final GRPO-Control Extension}
\label{sec:methodology-grpo-preparation}

The final project stage will introduce and evaluate a GRPO-based control variant. The role of the midway methodology is therefore to define the problem formulation, notation, baseline algorithms, implementation paths, and evaluation protocol before adding the final method. In the final report, the GRPO-control method should be evaluated under the same MDP formulation, seed structure, environment set, and evaluation metric as the PPO, SAC, and TD3 baselines.

This methodology therefore establishes the comparison framework rather than the final algorithmic contribution. The current implementation validates that the baseline infrastructure works end-to-end. The remaining methodological work is to formalize the GRPO-control update, motivate its components, and compare it against the completed baseline matrix.